\section{AI 收入从谁口袋里来：电子收入池与劳动力池}

前一章讨论各赛道，本章回答更根本的问题：这些 AI 公司的巨额收入从谁口袋里来？Freda 把互联网存量收入称作“电子收入/电子税”：线上广告、电商佣金和订阅，加起来约 4000 亿美元级别。这个池子不算小，但如果 OpenAI 一家公司要做到 2000 亿美元收入，光抢这个池子就意味着 Google、Meta、Amazon 等要让出大量份额。

\begin{figure}[H]
\centering
\includegraphics[width=0.94\textwidth]{electronic-revenue-vs-labor.png}
\caption{电子收入 vs 劳动力池：互联网存量不够大，真正大数在劳动和生产力。自制概念图，依据 01:09:51--01:14:11 对谈内容整理。}
\end{figure}

\begin{knowledgebox}{读图：为什么只抢广告不够大}
左侧电子收入包括广告、电商佣金和订阅，节目给出的量级约 4000 亿美元。右侧劳动力池则是美国 15 万亿美元级别的劳动成本。AI 如果只做“小 Google”，意义有限；如果进入客服、知识工作、AI for Science 和生产力提升，收入空间才会真正放大。
\end{knowledgebox}

\subsection{广告收入会不会因为 AI 变大}

Freda 不同意“AI 会让广告总额显著增加”的说法。美国广告收入过去二十年增长相对稳定，企业平均广告支出占收入比例也没有理由突然大幅提高，线上广告渗透率已经很高。因此，AI 抢广告更像存量重分配，而不是凭空扩张一个巨大新池子。

\begin{warningbox}{广告增长和广告份额要分开}
AI 可以改变谁拿到广告收入，但不一定显著扩大广告总盘子。ChatGPT、Google、TikTok、Amazon 和 Meta 的竞争，更多是份额和入口重分配。
\end{warningbox}

\subsection{劳动力池、生产力和 AI 税}

如果模型能力上来，真正的大数在劳动力成本。美国 GDP 约 30 万亿美元，劳动力成本约 15 万亿美元，单客服就是数千亿美元市场。AI 如果提升 10\% 全球生产力，可以带来巨大的 GDP 增量。Freda 用一个简单例子说明：100 个人产出 100 个商品，AI 后 80 个人产出 110 个更便宜更好的商品，real GDP 增长，公司利润和剩余员工收入可能提高，但被替代的 20 个人会经历痛苦转型。

因此节目也提到 AI 税或裁员税的可能性。它不是危言耸听，而是生产力革命常见的分配问题：总产出变大，不代表每个人都短期受益。

\begin{importantbox}{AI 收入的两个阶段}
早期，模型能力还没完全上来，AI 更容易抢电子收入：广告、订阅、电商佣金和软件预算。后期，如果 Agent 和 AI for Science 成熟，AI 才可能真正进入劳动力池和新增生产力池。
\end{importantbox}

\subsection{本章小结}

AI 收入的来源分两层。互联网存量收入能支撑早期商业化，但不够解释万亿级投入；真正的大数在劳动力成本、生产力提升和新价值创造。泡沫判断要看 AI 是否能从第一层走到第二层。

